\documentclass[aps,prd,twocolumn,superscriptaddress,nofootinbib]{revtex4-2}
\usepackage{amsmath,amssymb,amsthm}
\usepackage[colorlinks=true,linkcolor=blue,citecolor=blue,urlcolor=blue]{hyperref}

\newtheorem{theorem}{Theorem}
\newtheorem{proposition}[theorem]{Proposition}
\newtheorem{corollary}[theorem]{Corollary}
\theoremstyle{remark}
\newtheorem*{remark}{Remark}
\newcommand{\bW}{\beta_{\mathrm W}}
\newcommand{\Z}{\mathbb{Z}}
\newcommand{\R}{\mathbb{R}}

\begin{document}

\title{A gap is not a graviton: strong-coupling exclusions for lattice gravity in first-order form}

\author{Tristan Romanov}
\affiliation{Independent Researcher, San Francisco, CA, USA}
\email{romanov360@gmail.com}

\author{Brent Hartshorn}

\date{Draft, 31 October 2026}

\begin{abstract}
In Yang--Mills theory a mass gap is the goal; in gravity it is a diagnosis. We prove, in Lean~4 with Mathlib,
that a family of reflection-positive lattice theories whose transfer-matrix gap is bounded below in physical units
along a continuum tuning has no continuum limit with a massless mode, and that a massless mode makes the two-point
function decay more slowly than any exponential. Combined with the strong-coupling theorems of lattice Yang--Mills
for every compact simple group, this excludes the strong-coupling phase of Euclidean lattice gravity written as a
gauge theory of the Lorentz group $SO(4)$ (for $\bW<4/29$) or of the de Sitter group $SO(5)$ (for $\bW<15/58$, without
compensator) from being gravitational: these phases have a unique state, a positive transfer matrix (so no ghost
states), and a gap. We also prove in Lean that an antiscreening gravitational correction to the Yang--Mills
$\beta$-function preserves asymptotic freedom while a screening one produces an interacting ultraviolet fixed point;
that on the 600-cell a gauge-fixing forest can lower the single-link Dobrushin partner count only from 10 to 9, so
$\bW<4/9$ is optimal there; and that the MacDowell--Mansouri plaquette-pair action leaves 30 partners per link in
the same gauge. We determine exactly the spectrum of the two-dimensional causal-triangulation transfer matrix:
$\lambda_n=x_-^{2n+2}/g$, geometric with ratio $(1-s)/(1+s)$, $s=\sqrt{1-4g^2}$, so the lattice gap is of order the
spacing, as a continuum gap requires. Two-dimensional Yang--Mills on a Brownian disk has the boundary Wilson loop
$d_F(1+gL\sqrt{C_F})e^{-gL\sqrt{C_F}}$ in the quantum length $L$. On small SU(2) plaquette clusters read as spin
networks, every harmonic gap is $\sqrt{2\mu_{\min}}$ for the plaquette Laplacian $\partial\partial^{T}$, which proves
the $N$-plaquette strip formula; the largest eigenvalue of the triple-product volume operator approaches the
classical regular tetrahedron as $1-3/(4j)$; and the Barbero--Immirzi parameter is certified to 25 digits for the
three standard puncture countings. Nothing here identifies a gravitational phase or constructs a continuum theory.
\end{abstract}

\maketitle

\section{Introduction}

Four-dimensional lattice Yang--Mills theory has a unique state and a transfer-matrix gap at strong coupling for
every compact simple gauge group~\cite{4YM}, and the gap problem in physical units is reduced there to a single cube
estimate. Gravity is the next force to add~\cite{QGL}. The present note starts from one observation: the graviton
is massless, so in any approach where gravity is a lattice theory with a transfer matrix, a gap in physical units
is a proof that the phase is \emph{not} gravitational. The strong-coupling machinery of~\cite{4YM} produces such
gaps at explicit couplings, and some lattice formulations of gravity are gauge theories of compact groups to which
it applies as stated.

The formulations are first-order (Palatini) gravity, whose connection is an $SO(4)$ gauge field in Euclidean
signature, and MacDowell--Mansouri gravity~\cite{MM}, an $SO(5)$ gauge field broken to $SO(4)$ by a compensating
vector~\cite{SW}. Monte Carlo simulations of such models find a vanishing tetrad at strong coupling~\cite{CF}, and
causal and Euclidean dynamical triangulations have non-geometric phases~\cite{CDTreview}; none of this is proved.
The pure Lorentz sector of a Wilson plaquette action generates $R_{\mu\nu}R^{\mu\nu}$~\cite{Zubkov,Gwl}, the
curvature-squared sector of quadratic gravity, whose spin-2 ghost is conjectured to be confined at strong
coupling~\cite{HR,LQA}.

All Lean statements are in the repository \texttt{romanov360/quantum-gravity}, directory \texttt{QuantumGravity/},
built against Mathlib (Lean v4.33.0), without \texttt{sorry}, with the standard axioms only; numerical checks are in
\texttt{python/qg/}.

\section{A physical gap excludes a massless mode}

Let $\rho_a$ be the finite spectral measure, in lattice units, of the transfer matrix at spacing $a$ for a
spatially smeared gauge-invariant field, so that the temporal two-point function is
$S_a(t)=\int e^{-Et/a}\,d\rho_a(E)$. A gap $m$ in physical units along the tuning means $\rho_a((-\infty,ma))=0$.

\begin{theorem}[\texttt{physicalGap\_ratio\_le}, \texttt{limit\_ratio\_of\_physicalGap}]
For $0\le t_1\le t_2$, $S_a(t_2)\le e^{-m(t_2-t_1)}S_a(t_1)$ for every $a$, and if $Z_aS_a(t_i)\to L_i$ along any
filter for positive renormalisations $Z_a$, then $L_2\le e^{-m(t_2-t_1)}L_1$.
\end{theorem}

A massless particle in a reflection-positive theory gives the smeared field spectral weight at arbitrarily small
energies: a packet contains momenta $|p|<\varepsilon$ and $E=|p|$. We call a finite measure $\sigma$ on $[0,\infty)$
\emph{massless} if $\sigma([0,\varepsilon))>0$ for every $\varepsilon>0$.

\begin{theorem}[\texttt{massless\_not\_expDecay}]
If $\sigma$ is massless, then for no $\Delta>0$ and $C$ does $\int e^{-Et}d\sigma(E)\le Ce^{-\Delta t}$ hold for
all $t\ge1$.
\end{theorem}

\begin{proof}[Idea]
$\int e^{-Et}d\sigma\ge\sigma([0,\Delta/2))\,e^{-\Delta t/2}$, and $Ce^{-\Delta t/2}\to0$.
\end{proof}

\begin{theorem}[\texttt{no\_massless\_limit\_of\_physicalGap}]\label{thm:nograv}
If a family of lattice theories has a gap $m>0$ in physical units along the tuning and its renormalised two-point
functions converge at every $t>0$ to $\int e^{-Et}d\sigma(E)$ for a finite measure $\sigma$ on $[0,\infty)$, then
$\sigma$ is not massless.
\end{theorem}

The other alternative, a gap bounded below in \emph{lattice} units, is the no-go theorem of~\cite{4YM}: every
renormalised two-point function then vanishes in the limit. A graviton therefore needs $\Delta_{\rm lat}(a)/a\to0$.

\section{First-order lattice gravity at strong coupling}

The theorem of~\cite{4YM} on every compact simple group gives, in the forest gauge, a unique free-boundary state,
exponential clustering and a transfer-matrix gap for $\bW=N^2\beta<2N\,\mathrm{Ric}/29$, where $N$ is the dimension
of the representation in the Wilson action and $\mathrm{Ric}=c/4$ when the Killing form is $c$ times the trace form.
The constants are quadratic-form identities on a Cartan subalgebra, proved for all Cartan elements in
\texttt{GravityGroups.lean}:

\begin{center}
\begin{tabular}{llccc}
\hline
group (rep.) & role & $c$ & Ric & range of $\bW$\\
\hline
$SU(2)$ (2) & self-dual Spin(4) & 4 & 1 & $4/29$\\
$SO(4)$ (4) & Lorentz, tetrad index & 2 & $1/2$ & $4/29$\\
$SO(5)$ (5) & de Sitter (MM) & 3 & $3/4$ & $15/58$\\
$Sp(2)$ (4) & de Sitter, spinor & 6 & $3/2$ & $12/29$\\
\hline
\end{tabular}
\end{center}

$SO(4)$ is semisimple; its two factors have the same Ricci constant, which is all the Bakry--\'Emery step uses.

\begin{corollary}\label{cor:lorentz}
The Euclidean lattice Lorentz gauge theory with weight $\exp(4\beta\sum_p\mathrm{Tr}_4U_p)$, $U\in SO(4)$, has for
$\bW=16\beta<4/29$ a unique infinite-volume state, exponential clustering and a transfer-matrix gap. Its transfer
matrix is a positive operator (reflection positivity of the Wilson action), so its Hilbert space has no
negative-norm states. If the gap persists in physical units along a tuning, the limit has no massless mode in any
channel (Theorem~\ref{thm:nograv}).
\end{corollary}

The centre element $-1\in SO(4)$ acts as $-1$ on the vector representation, so vector Wilson loops have an area law
at strong coupling~\cite{OS}. The fundamental character coefficient (\texttt{lorentz\_strong.py}) is
$a_F(\bW)=\bW/16+O(\bW^3)$ (for $SU(2)$: $I_2(\bW)/I_1(\bW)=\bW/4+\dots$), so the leading-order string tension
at $\bW=4/29$ is $-\log a_F=4.75$ in lattice units.

\begin{remark}
In the frozen-tetrad model the curvature-squared sector is therefore ghost-free and gapped at strong coupling; this
is the fixed-spacing analogue of ghost confinement in quadratic gravity~\cite{HR,LQA}. Whether a continuum limit of
this lattice model is quadratic gravity is not addressed.
\end{remark}

For MacDowell--Mansouri gravity~\cite{MM} the action pairs each plaquette with the complementary plaquette at the
same site. In the forest gauge of~\cite{4YM} every free link then has exactly 30 partners with multiplicity (42
without gauge fixing; \texttt{MMPartners.lean}, by \texttt{decide}), against 12 for the Wilson action. The
Lipschitz constant of a pair term and the compensator remain; the programme is in
\texttt{docs/G2-pregeometric-phase.md}.

\section{Gravity and asymptotic freedom}

Above the Planck scale the Yang--Mills coupling runs as $dg/dt=-bg^3-fg$, with $f$ a gravitational coefficient
whose sign is scheme dependent but non-negative in symmetry-preserving computations~\cite{FLP}. In $u=1/g^2$,
$u(t)=(u_0+b/f)e^{2ft}-b/f$.

\begin{theorem}[\texttt{RunningWithGravity.lean}]
(i) $u$ solves $du/dt=2b+2fu$. (ii) If $f>0$, $b\ge0$, $u_0\ge0$ then $u(t)\ge u_0+2bt$: the coupling is never larger
than without gravity, and $g^2\to0$. (iii) If $f<0<b$, then $g^2\to-f/b>0$, an interacting ultraviolet fixed point.
(iv) If $f>0$, every $u_*\ge u_0$ is attained at $t_*=\log\frac{u_*+b/f}{u_0+b/f}/(2f)\ge0$.
\end{theorem}

Item (iv) is the statement of~\cite{4YM} that the intermediate coupling window is entered at a spacing-independent
scale, with the gravitational correction included.

\section{Yang--Mills on curved lattices}

The single-link Dobrushin count of~\cite{4YM} is local: on any 2-complex with a gauge-fixing forest $F$, a free link
$e$ has $P(e)=\sum_{p\ni e}(\#\text{free links of }p-1)$ partners and $SU(2)$ is unique for $\bW<4/\max_eP(e)$. The
search in \texttt{forest\_partners.py} reproduces $P=12$ on $\Z^4$ and $P=7$ for the period-3 forest on $\Z^3$, and
finds by exhaustion that period-2 forests on $\Z^3$ cannot go below 8.

\begin{proposition}[\texttt{cell600\_nine\_optimal}]
On the 600-cell (120 vertices, 720 edges, every edge in 5 triangles) every forest leaves some free link with at least
9 partners, and a forest with $\max P=9$ exists. The optimal single-link $SU(2)$ range is $\bW<4/9$.
\end{proposition}

\begin{proof}
$P(e)=10-n(e)$, where $n(e)$ is the number of fixed links sharing a triangle with $e$. A forest has at most 119
links, each a neighbour of 10 links, so $\sum_{e\notin F}n(e)\le1190<2\cdot601$; some free link has $n(e)\le1$.
A greedy forest attains 9.
\end{proof}

For closed triangulations a forest has $V-1$ links among $E\approx6V$, so gauge fixing helps little and the range is
governed by the maximal edge degree; a uniform statement over dynamical triangulations needs a degree cap.

\section{Two exactly solvable models}

\paragraph{YM$_2$ on a Brownian disk.} Two-dimensional Yang--Mills depends on the surface only through its
area~\cite{Sengupta,Levy}; the free Brownian disk of perimeter $L$ has area density
$L^3(2\pi A^5)^{-1/2}e^{-L^2/2A}$~\cite{BM}. Hence
$\mathbb E[e^{-sA}]=(1+L\sqrt{2s})e^{-L\sqrt{2s}}$ (an inverse-Gamma Laplace transform via $K_{3/2}$), and the
Haar-averaged boundary Wilson loop $d_Fe^{-g^2AC_F/2}$ becomes
\[
\langle\!\langle W_F\rangle\!\rangle_L=d_F\,(1+gL\sqrt{C_F})\,e^{-gL\sqrt{C_F}},
\]
a perimeter law in quantum length with tension $\sqrt{2\sigma_F}$. Since $\int Z(U,A)\,dU=1$ for every $A$,
quenched and annealed averages agree. On the sphere each representation shifts the cosmological constant by
$g^2C_\lambda/2$, and only the trivial one is non-analytic at $\mu=0$: the string susceptibility exponent $-1/2$ is
unchanged. Monte Carlo checks are in \texttt{ym2\_brownian.py}.

\paragraph{2D causal triangulations.} The one-step transfer matrix of 2D CDT is
$M(l_1,l_2)=g^{l_1+l_2-1}\binom{l_1+l_2-1}{l_1-1}$~\cite{AL}, $0<g<\tfrac12$.

\begin{theorem}[\texttt{CDT2D.lean}; proof in \texttt{docs/G3-cdt.md}]\label{thm:cdt}
Let $s=\sqrt{1-4g^2}$ and $x_-=(1-s)/(2g)$. The symmetrised transfer matrix $D^{-1/2}MD^{1/2}$, $D=\mathrm{diag}(l)$,
is a positive Hilbert--Schmidt operator on $\ell^2$ whose nonzero spectrum is $\lambda_n=x_-^{2n+2}/g$,
$n\ge0$. Hence $\lambda_{n+1}/\lambda_n=(1-s)/(1+s)$ and $\Delta_{\rm lat}=2\,\mathrm{artanh}\,s$; with
$g=\tfrac12e^{-a^2\Lambda}$, $\Delta_{\rm lat}(a)/a\to2\sqrt{2\Lambda}$.
\end{theorem}

\begin{proof}[Idea]
On $v_x=(x^l)_l$, $v_xM=h(x)\,v_{\varphi(x)}$ with $h=x/(1-gx)$ and the M\"obius map $\varphi=g/(1-gx)$ (proved in
Lean from the negative binomial series). Differentiating $n$ times at the fixed point $x_-\in(0,1)$, where
$\varphi'(x_-)=x_-^2$ and $h(x_-)=x_-^2/g$, makes $M$ triangular on $w_k=\partial^kv_x|_{x_-}$ with diagonal
$x_-^{2n+2}/g$. The $w_kD^{1/2}$ lie in $\ell^2$ and span a dense set (an analytic function on the unit disk with all
derivatives zero at $x_-$ vanishes), and a compact self-adjoint operator is diagonalised by them.
\end{proof}

This is a solvable model of quantum gravity in which the lattice gap scales with the spacing, as~\cite{4YM}
requires of a theory with a continuum gap; 2D gravity has no graviton, so Theorem~\ref{thm:nograv} is silent.

\section{Spin networks on small clusters}

Hartshorn~\cite{BHqg} reads the exactly solved SU(2) Hamiltonian clusters as loop-quantum-gravity kinematics: the
link Casimir $j(j+1)$ is a squared area in units of $8\pi\gamma\ell_P^2$, and the gauge-invariant states are spin
networks. Three of the open items listed there have the following answers.

\paragraph{Harmonic gaps of any cluster.} Abelianising at weak coupling, plaquette fluxes $B=\partial A$ are
oscillators with kinetic matrix $K=\partial\partial^T$ on plaquettes, and $2-W_p\simeq|B_p|^2/4$ gives
$\omega^2=\mu/2$, $\mu\in\mathrm{spec}\,K$ (three-fold). The lowest gauge-invariant level is two soft quanta:
$\Delta_{g\to0}=\sqrt{2\mu_{\min}}$, with $\mu_{\min}$ the least nonzero eigenvalue (zero modes are Bianchi
identities of closed surfaces). For planar clusters $K=4I-A$ with $A$ the adjacency of plaquettes sharing a link.
This reproduces $\sqrt8$, $\sqrt6$ and $2\sqrt{2-\cos\frac\pi4}$ for one, two and three plaquettes, explains why the
L-shape has the strip's gap (both graphs are $P_3$), and predicts $\Delta=2$ for the $2\times2$ block and
$2\sqrt{2-\sqrt3/2}$ for the T-shape. For the $N$-strip the eigenvectors are sines, and \texttt{ClusterModes.lean}
proves $\Delta_{g\to0}=2\sqrt{2-\cos\frac{\pi}{N+1}}$ for every $N$.

\paragraph{The volume of a four-valent vertex.} In the coupled basis $|k\rangle$ of $(j_1j_2)\to k$, the triple
product $Q=\epsilon_{abc}J_{1a}J_{2b}J_{3c}$ is $i$ times a tridiagonal matrix with
$a_k=\sqrt{\Delta(j_1,j_2,k)\Delta(j_3,j_4,k)}/(4\sqrt{4k^2-1})$,
$\Delta(a,b,k)=(a+b+k+1)(-a+b+k)(a-b+k)(a+b-k+1)$, which agrees with brute-force tensor products to $10^{-14}$.
Matching $V^2=\frac29|A_1\cdot(A_2\times A_3)|$ of a classical tetrahedron fixes
$V=\frac{\sqrt2}{3}(8\pi\gamma\ell_P^2)^{3/2}\sqrt{|Q|}$. For four equal spins the ratio of $\lambda_{\max}(Q)$ to
the regular tetrahedron's $\frac{4}{3\sqrt3}[j(j+1)]^{3/2}$, which~\cite{BHqg} found rising to 0.921 at $j=8$,
tends to 1: numerically $j(1-\text{ratio})\to3/4$ (0.74998 at $j=65536$), a semiclassical limit of the kind
derived by Bohr--Sommerfeld methods~\cite{BH}.

\paragraph{The Immirzi parameter.} Matching $S=A/4\ell_P^2$ to puncture counts gives $\gamma=\ln2/(\pi\sqrt3)$ for
spin-$\frac12$ punctures, and for all spins the roots of $1=\sum_j c_j e^{-2\pi\gamma\sqrt{j(j+1)}}$ with $c_j=2$
(Domagala--Lewandowski, Meissner~\cite{DL,Meissner}) or $c_j=2j+1$ (Ghosh--Mitra~\cite{GM}). Bisection with an
explicit geometric tail bound certifies $\gamma=0.2375329579659206643583771$ and $0.2740668582328299962167637$.

\section{Cosmological billiards}

Near a spacelike singularity, vacuum gravity becomes a billiard for the hyperbolic Kac--Moody algebra $AE_3$, and
eleven-dimensional supergravity one for $E_{10}$~\cite{DHN}. \texttt{Billiards.lean} checks by \texttt{decide} the
leading principal minors $2,3,-2$ of $AE_3$ and $2,3,\dots,10,-1$ of $E_{10}$ (one sign change: Lorentzian
signature), and \texttt{bkl\_billiard.py} that the BKL era lengths are the continued-fraction digits of the initial
$u$, with Gauss--Kuzmin statistics.

\section{What is not claimed}

No result here identifies a gravitational phase, constructs a continuum limit, or concerns Lorentzian signature.
The strong-coupling statements are at fixed lattice spacing; Theorem~\ref{thm:nograv} needs the gap along a tuning,
which is not proved for any of the models above. The open problems are listed in \texttt{OPEN.md}.

\begin{thebibliography}{99}
\bibitem{4YM} D.~Bhattacharjee, U.~Bhattacharya, J.~M.~Lockwood, T.~Romanov, \emph{Explicit lattice Yang--Mills mass
gap for every compact simple group reducing the continuum gap to one cube}, Zenodo
\href{https://doi.org/10.5281/zenodo.23444338}{10.5281/zenodo.23444338} (2026), code at github.com/creelie/4Y-M.
\bibitem{QGL} B.~Hartshorn, \emph{Quantum Gravity Lab}, github.com/brentharts/Quantum\_Gravity\_Lab (2026).
\bibitem{MM} S.~W.~MacDowell, F.~Mansouri, Phys.\ Rev.\ Lett.\ \textbf{38}, 739 (1977).
\bibitem{SW} K.~S.~Stelle, P.~C.~West, Phys.\ Rev.\ D \textbf{21}, 1466 (1980).
\bibitem{CF} S.~Catterall, D.~Ferrante, arXiv:0912.5525.
\bibitem{CDTreview} J.~Ambj\o rn, R.~Loll, arXiv:2401.09399.
\bibitem{Zubkov} M.~A.~Zubkov, arXiv:1308.3406.
\bibitem{Gwl} \emph{Gravitational Wilson loop in discrete quantum gravity}, arXiv:0907.2652.
\bibitem{HR} B.~Holdom, J.~Ren, Phys.\ Rev.\ D \textbf{93}, 124030 (2016).
\bibitem{LQA} R.~Liu, J.~Quintin, N.~Afshordi, Phys.\ Rev.\ Lett.\ \textbf{136}, 111501 (2026), arXiv:2510.18733.
\bibitem{OS} K.~Osterwalder, E.~Seiler, Ann.\ Phys.\ \textbf{110}, 440 (1978).
\bibitem{FLP} S.~Folkerts, D.~F.~Litim, J.~M.~Pawlowski, Phys.\ Lett.\ B \textbf{709}, 234 (2012), arXiv:1101.5552.
\bibitem{Sengupta} A.~Sengupta, Mem.\ Amer.\ Math.\ Soc.\ \textbf{126} (1997).
\bibitem{Levy} T.~L\'evy, Mem.\ Amer.\ Math.\ Soc.\ \textbf{166} (2003).
\bibitem{BM} J.~Bettinelli, G.~Miermont, Ann.\ Inst.\ H.\ Poincar\'e Probab.\ Statist.\ \textbf{53}, 1244 (2017).
\bibitem{AL} J.~Ambj\o rn, R.~Loll, Nucl.\ Phys.\ B \textbf{536}, 407 (1998).
\bibitem{BHqg} B.~Hartshorn, \emph{Spin networks on the SU(2) ladder: checked equations toward loop quantum
gravity}, github.com/brentharts/openai\_math\_yang\_mills (2026).
\bibitem{BH} E.~Bianchi, H.~M.~Haggard, Phys.\ Rev.\ Lett.\ \textbf{107}, 011301 (2011).
\bibitem{DL} M.~Domaga\l a, J.~Lewandowski, Class.\ Quant.\ Grav.\ \textbf{21}, 5233 (2004).
\bibitem{Meissner} K.~A.~Meissner, Class.\ Quant.\ Grav.\ \textbf{21}, 5245 (2004).
\bibitem{GM} A.~Ghosh, P.~Mitra, Phys.\ Lett.\ B \textbf{616}, 114 (2005).
\bibitem{DHN} T.~Damour, M.~Henneaux, H.~Nicolai, Class.\ Quant.\ Grav.\ \textbf{20}, R145 (2003).
\end{thebibliography}

\end{document}
